\section{Threat Model}
% \vspace{-5mm}
In this section, we define the threat model with respect to the attacker. We do this by describing the attacker's goal, background knowledge, and capabilities.

\textbf{Attacker's Goal.}

The attacker's primary objective is to manipulate the multi-agent system to derail the completion of benign tasks or trigger malicious actions aligned with their intent. Unlike single-agent settings, the attacker can achieve this indirectly by influencing inter-agent communication, disrupting coordination protocols, or exploiting specialized role assumptions, thereby causing system-wide cascading effects.


\textbf{Attacker's background knowledge.}
 
The attacker is assumed to know the roles and tools accessible to individual agents, but not the underlying LLM parameters such as alignment strategies, model parameters, and architectural details. Even this limited knowledge is sufficient to target weak links of the system, whose compromised outputs can propagate adversarially through the system.

\textbf{Attacker's capabilities.}

The attacker may (i) inject malicious content at the prompt or environment level, (ii) compromise one or more agents via adversarial system prompts, or (iii) add tools with malicious intent into the agent's toolkit. These capabilities enable attacks across three surfaces in the multi-agent system: prompt-level, agent-level, and environment-level.

\section{Attacks}
\subsection{Preliminaries}
We consider a multi-agent LLM system designed to handle user queries using collaborative agents. Let $q$ be the user query sampled from a distribution of queries $\pi_q$. Let $\mathcal{M}$ denote the multi-agent system consisting of $n$ agents $\{A_1, A_2, \dots, A_n\}$. Each agent $A_i$ is initialized with a system prompt $p^{\text{sys}}_i$ that defines its role, instructions, or behavioral constraints. $T_i = (\tau_i^1, \tau_i^2, \dots, \tau_i^n)$ denotes the list of tools available to agent $A_i$, where $T_i$ represents the set of agent-specific tools. An agent can invoke these tools to perform the user task. $O = (o_1, o_2, \dots, o_m)$ denotes the observations based on the actions taken by the agents. For a given query $q$ we aim to maximize:
\begin{equation}
  \mathbb{E}_{q \sim \pi_q} \left[ \mathds{1}\left( \mathcal{M}(q, O, \{T_i\}, \{p^{\text{sys}}_i\}) = a_b \right) \right]  
\end{equation}


where $a_b$ is the benign action and $\mathds{1}$ is an indicator function. A user aims to solve a target task $t$ consisting of an instruction, tools and data. The instruction corresponding to the target task is denoted using $q^t$.
% and the tools needed to solve the task are $T^t = (\tau_1^t, \tau_2^t,\dots, \tau_n^t)$.  

% \subsection{Attacks}
\label{sec:attacks}
\subsection{Prompt-level Attacks}
\subsubsection{Direct Prompt Injection (DPI)}
A DPI attack targets the multi-agent system by explicitly modifying the user query with a malicious instruction. In this attack, an injected instruction \( x^e \) is concatenated to the original target instruction \( q^t \), forming a new user query \( q^t \oplus x^e \). This injected instruction is crafted to override, or redirect the intended behavior of the system. Additionally, the adversary provides an attack-specific toolset \( T^e \), which is appended to the original tools \( T \). The goal of DPI is to manipulate the agents' behavior such that they follow the injected instruction instead of adhering to their original task. Formally, the goal is to maximize:
\begin{equation}
  \mathbb{E}_{q^t \sim \pi_q} \left[ \mathds{1}\left( \mathcal{M}(q^t \oplus x^e, O, \{T_i + T^e_i\}, \{p^{\text{sys}}_i\}) = a_{\text{m}} \right) \right]  
\end{equation}


where $a_m$ is the malicious action mentioned in the injected instruction. The attack is successful if the attack tool mentioned in the injected instruction is invoked.


\subsubsection{Impersonation}
An impersonation attack modifies the user query by appending a statement that falsely attributes the request to a trusted or authoritative figure. The modified query takes the form \( q^t \oplus x^{\text{auth}} \), where \( x^{\text{auth}} \) is a crafted phrase implying that the request originates from a legitimate or high-ranking source (e.g., “As requested by the admin...”). This tactic aims to make the query appear more valid or important, thereby increasing the likelihood that agents will comply with it, even if it contradicts their original role. The goal is to maximize:
\begin{equation}
    \mathbb{E}_{q^t \sim \pi_q} \left[ \mathds{1}\left( \mathcal{M}(q^t \oplus x^{\text{auth}}, O, \{T_i\}, \{p^{\text{sys}}_i\}) = a_{\text{m}} \right) \right]
\end{equation}

These adversarial behaviors would be rejected by well-aligned agents under normal, and unaltered conditions.

\subsection{Environment-level Attacks}
\subsubsection{Indirect Prompt Injection (IPI)}
An IPI attack manipulates the multi-agent system indirectly by introducing adversarial content into the environment or intermediary observations, rather than modifying the user query directly. In this setting, the user query remains unchanged as \( q^t \), but the attacker influences the observations \( O = (o_1, \dots, o_n)\) by injecting an instruction \( x^e \) at any step $i$ and appending to the attack toolset $T^e$ to $T$, resulting in altered observations.
% \( O = (o_1, \dots, o_i^e, \dots, o_n) \). 
These injected observations can originate from third-party tools or external sources accessed by agents during task execution. The goal of IPI is to mislead agents by feeding them altered or misleading context through compromised information. Formally, the goal is to maximize:
\begin{equation}
    \mathbb{E}_{q^t \sim \pi_q} \left[ \mathds{1}\left( \mathcal{M}(q^t, O \oplus x^e, \{T_i + T^e_i\}, \{p^{\text{sys}}_i\}) = a_{\text{m}} \right) \right]
\end{equation}


\subsection{Compromised Agents Attacks}
\subsubsection{Single Agent Compromise}
Single agent attacks occur when one agent in the multi-agent system is compromised, while the rest of the agents remain benign. Unlike prompt or environment based attacks, the adversarial influence arises solely from the malicious behavior of a single compromised agent. This setup highlights the system’s vulnerability to the weakest link: even one agent acting adversarially can mislead the overall decision-making process. Formally, this can be modeled by perturbing only the system prompt of the compromised agent as follows:
\begin{equation}
\small
\mathbb{E}_{q^t \sim \pi_q} \left[
    \mathds{1} \left(
        \mathcal{M}\big(
            q^t,\,
            O,\,
            \{T_i + T_i^e\},\,
            \{p^{\text{sys}}_1, \dots, p^{\text{sys}}_j + \delta_j, \dots, p^{\text{sys}}_N\}
        \big)
        = a_{\text{m}}
    \right)
\right]
\end{equation}
where $j$ denotes the index of the adversarial agent, whose system prompt $p^{\text{sys}}_j$ is modified with malicious instructions $\delta_j$, while all other agents remain unmodified.



\textbf{Byzantine agent.}  
A Byzantine agent directly produces inconsistent, or nonsensical outputs. This attack mode does not rely on persuasion or subtlety but rather on disrupting the system’s reasoning pipeline through contradictory, erroneous, or adversarially crafted outputs. Such an agent may provide factually incorrect answers, intentionally sabotage tool usage, or inject irrelevant noise into the communication. While Byzantine behavior is easier to detect than persuasive behavior, it can still reduce system robustness.



\subsubsection{Colluding Agents}
In a colluding agents attack, one or more agents within the multi-agent system are adversarial and deliberately coordinate to manipulate the system’s behavior toward an outcome desired by the attacker. These agents are initialized with adversarially modified system prompts of the form \( p^{\text{sys}}_i + \delta_i \), where \( \delta_i \) defines instructions encouraging the agents to cooperate toward an adversarial goal. The rest of the agents remain benign, but their outputs may be influenced or misled by the malicious agents through collaborative reasoning or message passing. Formally, the goal is to maximize:

\begin{equation}
\small
\mathbb{E}_{q^t \sim \pi_q} \left[
    \mathds{1} \left(
        \mathcal{M}\left(
            q^t,\,
            O,\,
            \{T_i + T_i^e\},\,
            \begin{aligned}
                &\{p^{\text{sys}}_i + \delta_i \mid i \in \mathcal{C}\} \\
                &\cup \{p^{\text{sys}}_i \mid i \notin \mathcal{C}\}
            \end{aligned}
        \right)
        = a_{\text{m}}
    \right)
\right]
\end{equation}

where, \( \mathcal{C} \subset \mathcal{M} \) denotes the set of colluding agents within the multi-agent system that intentionally cooperate to pursue a shared adversarial objective.

\subsubsection{Contradicting Agents}
In a contradicting agents attack, a subset of agents \( \mathcal{C} \subset \mathcal{M} \) which have similar functionalities, intentionally provide conflicting or misleading information to disrupt the overall system performance. Their goal is to derail the conversation, cause incomplete execution of the original target task, or generate adversarial responses by contradicting other agents. These agents modify their system prompts to \( p^{\text{sys}}_i + \delta_i \), where \( \delta_i \) defines the instructions to produce contradictory or disruptive behaviors. The goal is to maximize:

\begin{equation}
\small
\mathbb{E}_{q^t \sim \pi_q} \Biggl[
    \mathds{1} \Biggl(
        \mathcal{M}\Bigl(
            q^t,\, O,\, \{T_i\},\,
            \begin{aligned}[t]
            &\{p^{\text{sys}}_j + \delta_j,\, p^{\text{sys}}_k + \delta_k\} 
            \cup \{p^{\text{sys}}_i \mid i \notin \{j,k\}\}
            \end{aligned}
        \Bigr)
        = a_{\text{m}}\}
    \Biggr)
\Biggr]
\end{equation}



where, \( j,k \in \mathcal{C} \subset \mathcal{M} \) are two agents with similar functionalities that produce conflicting outputs. Here, $a_m$ can either correspond to an incomplete execution of the target task or an adversarial output.





\section{TAMAS Benchmark}

To evaluate the robustness of multi-agent systems we construct the \textbf{T}hreats and \textbf{A}ttacks in \textbf{M}ulti-\textbf{A}gent \textbf{S}ystems (\textbf{TAMAS}) benchmark. Following is an overview of the dataset and details the curation process. 

\subsection{Dataset Details}

\setlength{\intextsep}{1pt}
\begin{wraptable}{r}{0.55\textwidth}
\centering
\caption{TAMAS dataset statistics and representative instances, demonstrating coverage across agents, domains, attack types, and tasks.}
\resizebox{0.55\textwidth}{!}{%
\renewcommand{\arraystretch}{1.1}
% \begin{tabular}{lcl}
% \hline
% \multicolumn{1}{c}{\textbf{Category}} & \textbf{Value} & \multicolumn{1}{c}{\textbf{Example}} \\ \hline
% Scenarios & 5 & Education, Finance, News,\\ Healthcare and Legal. \\ \hline
% Attacks & 7 & \begin{tabular}[c]{@{}l@{}}Direct prompt injection, \\ colluding agents, etc\end{tabular} \\ \hline
% Datapoint per scenario & 50 & 10 datapoints for each attack \\ \hline
% \begin{tabular}[c]{@{}l@{}}Agents per Multi-Agent\\ System\end{tabular} & 4 & \begin{tabular}[c]{@{}l@{}}Legal Domain: Legal Research\\ Agent, Case Analysis Agent,\\ Document Drafting Agent, and\\ Compliance Check Agent.\end{tabular} \\ \hline
% Harmless tasks per scenario & 20 & \begin{tabular}[c]{@{}l@{}}Write an article about \\ advancements in AI for\\ healthcare and post it on\\ Linkedin and X.\end{tabular} \\ \hline
% \end{tabular}
\begin{tabular}{lcl}
\hline
\multicolumn{1}{c}{\textbf{Category}} & \textbf{Count} & \multicolumn{1}{c}{\textbf{Instance}} \\ \hline
\begin{tabular}[c]{@{}l@{}}Agents per\\ Multi-Agent System\end{tabular} & 4 & \begin{tabular}[c]{@{}l@{}}Legal Domain: Legal \\ Research Agent, Case \\ Analysis Agent, Document\\ Drafting Agent, and\\ Compliance Check Agent.\end{tabular} \\ \hline
Scenarios & 5 & \begin{tabular}[c]{@{}l@{}}Education, Finance, News,\\ Healthcare and Legal.\end{tabular} \\ \hline
Attacks & 6 & \begin{tabular}[c]{@{}l@{}}Direct prompt injection, \\ Indirect prompt injection,\\ Impersonation, Byzantine \\ agent, Colluding agents, and\\ Contradicting agents\end{tabular} \\ \hline
Harmless tasks per scenario & 20 & \begin{tabular}[c]{@{}l@{}}Write an article about \\ advancements in AI for\\ healthcare and post it on\\ Linkedin and X.\end{tabular} \\ \hline
Datapoint per scenario & 60 & \begin{tabular}[c]{@{}l@{}}10 datapoints for each\\  attack\end{tabular} \\ \hline
\end{tabular}
}

\label{tab:dataset_details}
\end{wraptable}


% \begin{wraptable}{r}{0.55\textwidth}
% \centering
% \resizebox{0.55\textwidth}{!}{%
% \renewcommand{\arraystretch}{1.4}
% \begin{tabular}{lcl}
% \hline
% \multicolumn{1}{c}{\textbf{Category}} & \textbf{Value} & \multicolumn{1}{c}{\textbf{Example}} \\ \hline
% Scenarios & 5 & Education, Finance, etc. \\
% Attacks & 7 & \begin{tabular}[c]{@{}l@{}}Direct prompt injection, \\ colluding agents, etc\end{tabular} \\
% Datapoint per scenario & 50 & 10 datapoints for each attack \\
% \begin{tabular}[c]{@{}l@{}}Agents per Multi-Agent\\ System\end{tabular} & 4 & \begin{tabular}[c]{@{}l@{}}Legal Domain: Legal Research\\ Agent, Case Analysis Agent,\\ Document Drafting Agent, etc.\end{tabular} \\
% Harmless tasks & 100 & \begin{tabular}[c]{@{}l@{}}Write an article about \\ advancements in AI for\\ healthcare and post it on\\ Linkedin and X.\end{tabular} \\ \hline
% \end{tabular}
% }
% \caption{Overview of dataset details and representative examples for each category.}
% \label{tab:dataset_details}
% \end{wraptable}




\textbf{Scenarios}: We construct a dataset spanning five real-world domain scenarios: News, Education, Finance, Healthcare, and Legal. These domains were selected to reflect diverse, high-stakes applications where LLM-based multi-agent systems are likely to be deployed and where safety and robustness are critical. Each scenario is built around a single multi-agent system composed of four distinct agents, each with specialized and diverse functionalities. For each scenario, we include ten adversarial examples per attack mentioned in Section \ref{sec:attacks}. Each datapoint consists of a multi-step task involving atleast two-three agents to capture the dynamics and inter-agent interactions in a multi-agent system. 

\textbf{Harmless instructions}: To assess the utility of the system we also include 20 harmless instructions per scenario.  These instructions reflect typical, non-adversarial tasks that a multi-agent system might encounter in the real world.  

\textbf{Synthetic Tools}: The actions performed by agents are enabled through a set of tools that each agent can access. These tools allow individual agents to perform tasks to fulfill the user query. The tools available to each agent depend on the domain and the role of the agent in the multi-agent system. We include two types of tools: (i) Normal Tools, which support standard execution of normal tasks (ii) Attack tools, which simulate malicious behavior.\\

All data, attack implementations, and evaluation scripts will be publicly released to support reproducibility and future research. Table \ref{tab:dataset_details} presents an overview of the benchmark, and additional details are included in Appendix \ref{App:dataset}.


\subsection{Agent Interaction Configurations}
We evaluate our dataset across three diverse agent interaction configurations to understand how these setups affect the susceptibility to adversarial attacks. We consider the following configurations from the Autogen \citep{wu2023autogen} and \href{https://docs.crewai.com/}{CrewAI} frameworks  for our study:

\textbf{Central Orchestrator}: In a centralized coordination paradigm, a lead orchestrator manages the overall workflow of the multi-agent system. The orchestrator is responsible for high-level planning, delegation of subtasks, and monitoring progress toward task completion. It begins by analyzing the user query to extract key requirements and formulate a structured plan. Each step of the plan is then assigned to the most suitable agent, while the orchestrator maintains a record of progress to ensure that subtasks are executed in the intended sequence. Once subtasks are completed, the orchestrator updates its progress tracker and continues to the next stage. By routing all decisions and interactions through a central entity, this design enforces structured control, accountability, and oversight across the system. We evaluate Magentic-One from Autogen and the centralized configuration from CrewAI.

\textbf{Sequential}: Employs a decentralized coordination strategy in which agents take turns contributing to the task in a fixed, cyclic order. After an agent completes its turn, control is passed on to the next agent in the sequence. This configuration employs equal participation, but lacks centralized planning and oversight of the tasks. We evaluate the Round Robin workflow of AutoGen framework and sequential configuration from CrewAI. 

\textbf{Collaborative}: Employs a dynamic coordination where the agents take turns contributing to the task at hand based on handoff decisions. In contrast to Round Robin configuration where the sequence of the agents was fixed, the agents in a Swarm configuration select the next agent through a handoff message by the current agent. This makes the configuration decentralized, yet adaptive in turn taking. All agents share a common message context, ensuring a consistent view of the task. Each agent is capable of signaling a handoff to another agent, enabling more flexible and context-sensitive coordination. We evaluate Swarm from Autogen. CrewAI does not provide an equivalent configuration. A summary of the key features of each configuration is provided in Table \ref{tab:differences} in Appendix \ref{sec:app_config}.

\subsection{Evaluation Setup}
\textbf{Models :}  We evaluate performance across ten LLMs: (i) GPT-4 \citep{openai2024gpt4technicalreport} (ii) GPT-4o \citep{openai2024gpt4ocard} (iii) GPT-4o-mini \citep{openai2024gpt4ocard} (iv) Gemini-2.0-Flash (v) Gemini-2.0-Flash-lite (vi) Deepseek-R1-32B \citep{deepseekai2025deepseekr1incentivizingreasoningcapability} (vii) Mixtral-8x7B-Instruct \citep{jiang2024mixtralexperts} (viii) Qwen3-32B \citep{yang2025qwen3technicalreport} (ix) Qwen3-8B \citep{yang2025qwen3technicalreport} and (x) Llama-3.1-8B-Instruct. Refer Appendix \ref{Appendix:models} for more details.


% \subsection{Evaluation Metrics}
\textbf{Evaluation Metrics :} Our evaluation considers 3 criteria: (i) \textit{Attack success}: whether the attack succeeds, reflecting the safety and robustness of the multi-agent system (ii) \textit{Task Completion in Benign Conditions}: whether the system can successfully complete a given task under no-attack conditions (iii) \textit{Robustness}: the system's ability to prevent attacks while maintaining task performance.

To assess robustness, we adopt the Agent Risk Assessment (ARIA) framework \citep{tur2025safearenaevaluatingsafetyautonomous}, which categorizes system responses into four outcomes: ARIA-1 represents immediate refusal, indicating safe behavior; ARIA-2 denotes delayed refusal, where the system hesitates before rejecting the task; ARIA-3 reflects an intent to complete the task but a failure to do so; and ARIA-4 indicates a successful attack, representing unsafe behavior. Due to the semantic complexity of the logs, we adopt an LLM-as-judge setup along with tool invocation checks to ensure scalable evaluation.


To assess the utility of the backbone LLM and the multi-agent system, we evaluate the system's ability to successfully complete harmless tasks in a no-attack setting. This is captured through the Performance under No Attack (PNA) metric. Lastly, we define Effective Robustness Score (ERS) as a composite metric that captures both safety and utility. It reflects the system's ability to remain functional and secure under adversarial conditions. Refer Appendix \ref{App:metrics} for additional details on metrics and human verification of LLM-as-judge



